\documentclass[12pt]{article}
\usepackage[utf8]{inputenc}
\usepackage[T1]{fontenc}
\usepackage{hyperref}
\title{Iamalele Language Resources for K-12 Curriculum Development}
\author{Compiled by Claude Sonnet 5.5 (claude-sonnet-5-5) for Andr\'as Kornai}
\date{}
\begin{document}
\maketitle

\section{Introduction}
Iamalele (also Yamalele; ISO 639-3 yml) is an Austronesian language of Fergusson Island, Milne Bay Province, Papua New Guinea. Wikipedia gives about 2,800 speakers, based on a 1987 count. This report rests on a single web search. Very little material was found.

\section{Primary Educational Resources}
No primers, textbooks or school materials in Iamalele were found in the search results.

\section{Linguistic and Reference Materials}
\begin{itemize}
\item Yamalele language (Wikipedia): \url{https://en.wikipedia.org/wiki/Yamalele_language}
\item Numeral survey entry for Iamalele (MPG), which records only three numerals: \url{https://lingweb.eva.mpg.de/channumerals/Iamalele.htm}
\item Glosbe language pages (online dictionary aggregator, quality unverified): \url{https://glosbe.com/yml/nmx}
\end{itemize}

\section{Religious and Community Resources}
\begin{itemize}
\item Joshua Project profile, Yamalele and Maiodom in Papua New Guinea: \url{https://joshuaproject.net/people_groups/15969}
\end{itemize}

\section{Digital and Audio Resources}
No audio, video or corpus resources were found.

\section{Curriculum Development Context}
The search did not confirm Iamalele-medium instruction in any school. The MPG numeral note suggests that speakers often use English numbers beyond a few small ones. Tok Pisin and English are used for wider contact. Any curriculum would need new documentation work with the community.
\end{document}